

\section*{Exercise 3.3}
Suppose $A\subset\mathbb{R}^n$ and let $(f_i)_{i=0}^\infty$ with $f_i:A\to\mathbb{R}^m$.
Classify whether each statement is equivalent to (i) pointwise convergence $f_i\to f$,
(ii) uniform convergence $f_i\to f$ uniformly, or (iii) neither.

\begin{enumerate}
\renewcommand\labelenumi{(\alph{enumi})}

\item Given $\varepsilon>0$, there exists $N\in\mathbb{N}$ such that for all $i\ge N$:
\[
\sup_{x\in A}\|f_i(x)-f(x)\|<\varepsilon.
\]
\textit{Classification: (ii) uniform convergence.} This is exactly the definition via the sup norm.

\item For all $x\in A$, for all $i\in\mathbb{N}$ there exists $\varepsilon>0$ such that
\[
\|f_i(x)-f(x)\|<\varepsilon.
\]
\textit{Classification: (iii) neither.} The statement is always true (take, e.g., $\varepsilon=\|f_i(x)-f(x)\|+1$), so it does not characterize convergence.

\item There exists $N\in\mathbb{N}$ such that for all $x\in A$, $i\ge N$, $\varepsilon>0$:
\[
\|f_i(x)-f(x)\|<\varepsilon.
\]
\textit{Classification: (iii) neither.} Since the left-hand side is nonnegative, the inequality for \emph{all} $\varepsilon>0$ forces $\|f_i(x)-f(x)\|=0$; i.e.\ there is $N$ with $f_i\equiv f$ on $A$ for all $i\ge N$. This condition is much stronger than uniform convergence and is not equivalent to it in general.

\item[\textit{*}d)] $\forall\varepsilon>0\ \forall x\in A\ \exists N\in\mathbb{N}\ \forall i\ge N:\ \|f_i(x)-f(x)\|<\varepsilon$.\\
\textit{Classification: (i) pointwise convergence.} The quantifier order matches the definition; $N$ may depend on $x$ and $\varepsilon$.

\item[\textit{*}e)] $\forall\varepsilon>0\ \exists N\in\mathbb{N}\ \forall x\in A\ \forall i\ge N:\ \|f_i(x)-f(x)\|<\varepsilon$.\\
\textit{Classification: (ii) uniform convergence.} Here $N$ works simultaneously for all $x\in A$.

\item[\textit{*}f)] $\forall\varepsilon>0\ \exists N\in\mathbb{N}\ \forall x\in A:\ (\,i\ge N\ \Leftrightarrow\ \|f_i(x)-f(x)\|<\varepsilon\,)$.\\
\textit{Classification: (iii) neither.} This is far stronger than uniform convergence (it also demands that for $i<N$ the error is $\ge\varepsilon$ for every $x$), and it need not hold even when $f_i\to f$ uniformly.
\end{enumerate}


\section*{Exercise 3.4}
Suppose $A\subset\mathbb{R}^n$ and let $(f_i)_{i=0}^\infty$ with $f_i:A\to\mathbb{R}^m$ be a sequence of functions.
Assume $f_i\to f$ \emph{uniformly} on $A$. Show that if $B\subset A$, then $f_i\to f$ uniformly on $B$.

\textit{Solution.} By uniform convergence on $A$, for every $\varepsilon>0$ there exists $N\in\mathbb{N}$ such that
for all $i\ge N$ and all $x\in A$ we have $\|f_i(x)-f(x)\|<\varepsilon$. Since $B\subset A$, this inequality holds
for all $x\in B$ as well, with the \emph{same} $N$. Hence $f_i\to f$ uniformly on $B$. Equivalently,
\[
\sup_{x\in B}\|f_i(x)-f(x)\|\le \sup_{x\in A}\|f_i(x)-f(x)\|\xrightarrow{i\to\infty} 0.
\]

